\section{Implementation and reporting conventions}\label{app:implementation}

This supplement records the implementation details and result grids that support the article. Every numerical cell is a direct transcription or a compact rearrangement of the finalized analysis underlying the article.

\subsection{Sample, timing, and portfolio design}

Table~\ref{tab:app_design_summary} collects the choices that apply across the learned models. The chart CNN, CNN1D, and multimodal models use the same CRSP source panel, target definition, sample split, and downstream portfolio rules. A method enters a formation-date cross section when its required input history and subsequent holding-period return are available.

\begin{table}[htbp]
\centering
\caption{Empirical design summary}
\label{tab:app_design_summary}
\small
\begin{tabularx}{\textwidth}{@{}p{0.25\textwidth}L@{}}
\toprule
Component & Convention \\
\midrule
Source universe & CRSP ordinary common shares with share codes 10 and 11 on NYSE, AMEX, or Nasdaq. Prices must be positive, and lagged market capitalization must be observed for value-weighted portfolios. \\
Development sample & 1993--2000. Stock-date observations are randomly divided into 70\% training and 30\% validation sets using a fixed seed. \\
Evaluation samples & Full period 2001--2024, baseline period 2001--2019, and extension period 2020--2024. \\
Input and target grid & Input length $I\in\{5,20,60\}$ trading days and future return horizon $R\in\{5,20,60\}$ trading days. The target equals one when the cumulative return over dates $t+1$ through $t+R$ is positive. \\
Portfolio frequency & Non-overlapping weekly, monthly, and quarterly formation dates for $R=5$, 20, and 60. \\
Portfolio rule & Buy the highest-score decile and short the lowest-score decile. Each leg has one dollar of notional. Results use equal- or lagged-market-capitalization weights within each leg. \\
Annualization & Mean returns use factors 52, 12, and 4. Standard deviations use the square roots of the same factors. The Sharpe ratio is annualized mean divided by annualized standard deviation. \\
Transaction costs & A one-way rate of 10 basis points applies above the contemporaneous NYSE 80th-percentile market-capitalization breakpoint and 20 basis points otherwise. The deduction uses full absolute traded notional in both legs. \\
Turnover & Conventional half-turnover is computed separately for the long and short legs, summed, and expressed as a monthly equivalent. \\
\bottomrule
\end{tabularx}
\end{table}

\subsection{Representations and network architectures}

The chart input contains binary open-high-low-close bars, a moving-average line with length $I$, and a separate volume panel. Prices and the moving average are jointly min--max scaled within the stock-window, and volume is scaled separately. The image-scaled sequence retains the six continuous open, high, low, close, moving-average, and volume channels after the corresponding within-window scaling. CNN1D applies convolution along the ordered trading-day dimension.

\begin{table}[htbp]
\centering
\caption{Architecture by input length}
\label{tab:app_architecture}
\small
\setlength{\tabcolsep}{5pt}
\begin{tabular}{@{}cclcl@{}}
\toprule
$I$ & Chart input & Chart CNN blocks and channels & Sequence input & CNN1D blocks and channels \\
\midrule
5  & $32\times15$  & 2 blocks, 64 and 128 & $6\times5$  & 1 block, 64 \\
20 & $64\times60$  & 3 blocks, 64, 128, and 256 & $6\times20$ & 2 blocks, 64 and 128 \\
60 & $96\times180$ & 4 blocks, 64, 128, 256, and 512 & $6\times60$ & 3 blocks, 64, 128, and 256 \\
\bottomrule
\end{tabular}
\begin{minipage}{\textwidth}
\footnotesize\textit{Note.} Every chart block applies a $5\times3$ convolution, batch normalization, Leaky ReLU activation, and max pooling. Every sequence block applies the corresponding operations with a one-dimensional filter of length three, stride one, and padding one. The final feature map is flattened, dropout is applied, and a fully connected layer produces two return-sign logits. The multimodal model pairs the chart and sequence branches for the same $(i,t,I,R)$ observation. Its primary fusion head concatenates the two branch feature vectors and applies a linear classifier.
\end{minipage}
\end{table}

The scaling exercise also estimates a logistic classifier on the image-scaled, cumulative-return-scaled, and devolatized-return-scaled sequences. Scheme-specific sequence channels are standardized with moments estimated in the development sample and then held fixed during evaluation.

\begin{table}[htbp]
\centering
\caption{Estimation settings for learned models}
\label{tab:app_training}
\small
\begin{tabularx}{0.86\textwidth}{@{}lL@{}}
\toprule
Setting & Value \\
\midrule
Objective & Two-class cross-entropy loss \\
Optimizer & Adam \\
Learning rate & $10^{-4}$ \\
Batch size & 128 observations \\
Maximum epochs & 50 \\
Checkpoint rule & Lowest validation loss \\
Early stopping & Two consecutive epochs without validation-loss improvement \\
Dropout probability & 0.50 after the flattened feature representation \\
Initialization & Xavier initialization for convolutional and fully connected weights \\
Portfolio score & Softmax probability assigned to the positive-return class \\
\bottomrule
\end{tabularx}
\end{table}

\subsection{Checkpoint and ensemble conventions}\label{app:run_conventions}

The main empirical analysis uses one selected checkpoint for every learned specification. The chart-CNN seed sensitivity exercise separately compares one- and five-member forecast averages for the two short-input weekly specifications. Table~\ref{tab:app_run_conventions} keeps the main results and this sensitivity check distinct.

\begin{table}[htbp]
\centering
\caption{Run convention by result family}
\label{tab:app_run_conventions}
\small
\begin{tabularx}{\textwidth}{@{}>{\raggedright\arraybackslash}p{0.24\textwidth}>{\raggedright\arraybackslash}p{0.23\textwidth}L@{}}
\toprule
Result family & Forecast aggregation & Scope \\
\midrule
Weekly mechanism grid & One selected checkpoint & All chart CNN, CNN1D, and multimodal $I5/R5$, $I20/R5$, and $I60/R5$ cells \\
Chart grid & One selected checkpoint & All nine $(I,R)$ cells \\
CNN1D grid & One selected checkpoint & All nine $(I,R)$ cells \\
Multimodal grid & One selected checkpoint & All nine feature-concatenation cells \\
Seed sensitivity & Matched $M=1$ and $M=5$ results & Chart CNN $I5/R5$ and $I20/R5$ \\
Late fusion & Prediction-level combination & Pretrained chart and CNN1D probabilities or logits are averaged after branch prediction \\
\bottomrule
\end{tabularx}
\end{table}

\FloatBarrier

\section{Complete representation decomposition}\label{app:scaling_grid}

Table~\ref{tab:app_scaling_grid} reports the full scaling exercise from the canonical analysis. It extends the article's weekly main grid to monthly and quarterly return horizons while preserving both portfolio weighting rules.

\begin{table}[htbp]
\centering
\caption{Gross Sharpe ratios across sequence scaling schemes}
\label{tab:app_scaling_grid}
\scriptsize
\setlength{\tabcolsep}{4.7pt}
\begin{tabular}{@{}lrrrrrr@{}}
\toprule
& \multicolumn{2}{c}{5-day input} & \multicolumn{2}{c}{20-day input} & \multicolumn{2}{c}{60-day input} \\
\cmidrule(lr){2-3}\cmidrule(lr){4-5}\cmidrule(l){6-7}
Model & EW & VW & EW & VW & EW & VW \\
\midrule
\multicolumn{7}{@{}l}{\textit{Panel A. 5-day return horizon}} \\
Chart CNN & 5.88 & 1.39 & 6.52 & 1.44 & 3.72 & 1.18 \\
Logistic, local image scale & 4.88 & 1.25 & 5.22 & 1.09 & 5.44 & 1.12 \\
Logistic, cumulative scale & 0.18 & $-0.04$ & $-0.41$ & $-0.51$ & $-0.26$ & $-0.53$ \\
Logistic, devolatized scale & 2.80 & 0.79 & 2.73 & 0.48 & 2.70 & 0.66 \\
CNN1D, local image scale & 6.24 & 1.16 & 6.55 & 1.12 & 5.89 & 2.06 \\
CNN1D, cumulative scale & 4.88 & 1.47 & 5.12 & 1.41 & 4.66 & 1.15 \\
CNN1D, devolatized scale & 4.39 & 0.73 & 5.26 & 1.04 & 4.82 & 1.28 \\
\addlinespace[3pt]
\multicolumn{7}{@{}l}{\textit{Panel B. 20-day return horizon}} \\
Chart CNN & 1.35 & 0.10 & 2.26 & 0.49 & 0.57 & $-0.29$ \\
Logistic, local image scale & 2.03 & 0.56 & 2.37 & 0.41 & 1.59 & 0.60 \\
Logistic, cumulative scale & 0.02 & 0.37 & $-0.29$ & $-0.08$ & $-0.30$ & $-0.64$ \\
Logistic, devolatized scale & 1.27 & 0.59 & 1.02 & 0.28 & 0.98 & 0.12 \\
CNN1D, local image scale & 2.18 & 0.41 & 2.15 & 0.41 & 1.74 & 0.51 \\
CNN1D, cumulative scale & 0.94 & 0.92 & 0.97 & 0.56 & 1.45 & 0.63 \\
CNN1D, devolatized scale & 1.17 & 0.16 & 1.73 & 0.53 & 1.22 & 0.30 \\
\addlinespace[3pt]
\multicolumn{7}{@{}l}{\textit{Panel C. 60-day return horizon}} \\
Chart CNN & 0.23 & $-0.18$ & 0.34 & 0.21 & 0.55 & 0.45 \\
Logistic, local image scale & 0.97 & 0.04 & 0.94 & 0.21 & 0.55 & 0.21 \\
Logistic, cumulative scale & 0.37 & 0.11 & $-0.02$ & $-0.06$ & $-0.30$ & $-0.54$ \\
Logistic, devolatized scale & 0.66 & 0.08 & 0.61 & 0.37 & 0.40 & 0.17 \\
CNN1D, local image scale & 0.99 & 0.28 & 0.46 & 0.15 & 0.46 & 0.25 \\
CNN1D, cumulative scale & 0.26 & 0.33 & 0.13 & 0.16 & 0.34 & 0.21 \\
CNN1D, devolatized scale & 0.83 & 0.13 & 0.65 & 0.34 & 0.45 & 0.21 \\
\bottomrule
\end{tabular}
\begin{minipage}{\textwidth}
\footnotesize\textit{Note.} Entries are annualized gross Sharpe ratios of high-minus-low decile portfolios over 2001--2024. EW and VW denote equal- and value-weighted legs. Local image scale applies the chart's within-window high--low normalization to the continuous sequence. Cumulative scale anchors the path to the first close, and devolatized scale expresses changes relative to lagged volatility.
\end{minipage}
\end{table}

\FloatBarrier

\section{Complete learned-model grids}\label{app:learned_grids}

Tables~\ref{tab:app_chart_grid}--\ref{tab:app_mmd_grid} condense the baseline and extension results into complete $3\times3$ grids. Gross SR and Net SR denote annualized Sharpe ratios before and after the common transaction-cost deduction. This layout keeps horizon, sample period, and weighting visible without repeating the return and standard-deviation columns already summarized in the article's focused economic tests.

\begin{table}[htbp]
\centering
\caption{Chart CNN Sharpe-ratio grid}
\label{tab:app_chart_grid}
\scriptsize
\setlength{\tabcolsep}{4.2pt}
\begin{tabular}{@{}lrrrrrrrr@{}}
\toprule
& \multicolumn{4}{c}{Equal-weighted} & \multicolumn{4}{c}{Value-weighted} \\
\cmidrule(lr){2-5}\cmidrule(l){6-9}
& \multicolumn{2}{c}{2001--2019} & \multicolumn{2}{c}{2020--2024}
& \multicolumn{2}{c}{2001--2019} & \multicolumn{2}{c}{2020--2024} \\
\cmidrule(lr){2-3}\cmidrule(lr){4-5}\cmidrule(lr){6-7}\cmidrule(l){8-9}
Spec. & Gross SR & Net SR & Gross SR & Net SR & Gross SR & Net SR & Gross SR & Net SR \\
\midrule
$I5/R5$       & 7.00 & 4.28 & 3.30 & 1.01 & 1.77 & 0.00 & 0.57 & $-0.47$ \\
$I5/R20$      & 1.98 & 1.05 & 0.21 & $-0.37$ & 0.26 & $-0.23$ & $-0.16$ & $-0.41$ \\
$I5/R60$      & 0.14 & $-0.10$ & 0.56 & 0.28 & $-0.15$ & $-0.28$ & 0.14 & $-0.03$ \\
$I20/R5$      & 7.64 & 4.78 & 3.87 & 1.35 & 1.71 & $-0.01$ & 0.70 & $-0.62$ \\
$I20/R20$     & 2.53 & 1.51 & 1.76 & 1.30 & 0.26 & $-0.31$ & 0.67 & 0.54 \\
$I20/R60$     & 0.35 & 0.16 & 0.02 & $-0.15$ & 0.25 & 0.08 & 0.12 & $-0.01$ \\
$I60/R5$      & 5.32 & 2.26 & 1.18 & $-0.61$ & 1.27 & $-0.68$ & 1.13 & $-0.45$ \\
$I60/R20$     & 1.28 & 0.31 & $-0.51$ & $-0.98$ & $-0.41$ & $-0.97$ & $-0.22$ & $-0.51$ \\
$I60/R60$     & 0.25 & 0.07 & 0.46 & 0.36 & 0.05 & $-0.01$ & 0.74 & 0.53 \\
\bottomrule
\end{tabular}
\begin{minipage}{\textwidth}
\footnotesize\textit{Note.} Each cell reports one selected chart-CNN specification from the canonical analysis. Returns are formed at the frequency associated with $R$ and are annualized using 52, 12, or 4 observations per year.
\end{minipage}
\end{table}

\begin{table}[htbp]
\centering
\caption{CNN1D Sharpe-ratio grid}
\label{tab:app_cnn1d_grid}
\scriptsize
\setlength{\tabcolsep}{4.2pt}
\begin{tabular}{@{}lrrrrrrrr@{}}
\toprule
& \multicolumn{4}{c}{Equal-weighted} & \multicolumn{4}{c}{Value-weighted} \\
\cmidrule(lr){2-5}\cmidrule(l){6-9}
& \multicolumn{2}{c}{2001--2019} & \multicolumn{2}{c}{2020--2024}
& \multicolumn{2}{c}{2001--2019} & \multicolumn{2}{c}{2020--2024} \\
\cmidrule(lr){2-3}\cmidrule(lr){4-5}\cmidrule(lr){6-7}\cmidrule(l){8-9}
Spec. & Gross SR & Net SR & Gross SR & Net SR & Gross SR & Net SR & Gross SR & Net SR \\
\midrule
$I5/R5$   & 7.45 & 4.50 & 3.47 & 1.02 & 1.37 & $-0.21$ & 0.54 & $-0.49$ \\
$I5/R20$  & 2.78 & 1.78 & 1.26 & 0.67 & 0.46 & 0.00 & 0.52 & 0.19 \\
$I5/R60$  & 1.08 & 0.83 & 0.73 & 0.51 & 0.44 & 0.28 & $-0.07$ & $-0.17$ \\
$I20/R5$  & 8.09 & 5.28 & 3.34 & 1.18 & 1.49 & $-0.35$ & 0.28 & $-0.77$ \\
$I20/R20$ & 2.63 & 1.92 & 1.30 & 0.85 & 0.38 & $-0.02$ & 0.70 & 0.25 \\
$I20/R60$ & 0.49 & 0.32 & 0.37 & 0.24 & 0.32 & 0.19 & $-0.31$ & $-0.45$ \\
$I60/R5$  & 7.88 & 5.58 & 2.29 & 0.78 & 2.61 & 0.80 & 0.92 & $-0.29$ \\
$I60/R20$ & 2.07 & 1.59 & 1.21 & 0.92 & 0.60 & 0.26 & 0.29 & $-0.02$ \\
$I60/R60$ & 0.49 & 0.37 & 0.44 & 0.32 & 0.36 & 0.23 & 0.06 & $-0.03$ \\
\bottomrule
\end{tabular}
\begin{minipage}{\textwidth}
\footnotesize\textit{Note.} Every cell uses one selected checkpoint of the image-scaled CNN1D model. Gross and net Sharpe ratios follow the portfolio and transaction-cost conventions in Table~\ref{tab:app_design_summary}.
\end{minipage}
\end{table}

\begin{table}[htbp]
\centering
\caption{Multimodal Sharpe-ratio grid}
\label{tab:app_mmd_grid}
\scriptsize
\setlength{\tabcolsep}{4.2pt}
\begin{tabular}{@{}lrrrrrrrr@{}}
\toprule
& \multicolumn{4}{c}{Equal-weighted} & \multicolumn{4}{c}{Value-weighted} \\
\cmidrule(lr){2-5}\cmidrule(l){6-9}
& \multicolumn{2}{c}{2001--2019} & \multicolumn{2}{c}{2020--2024}
& \multicolumn{2}{c}{2001--2019} & \multicolumn{2}{c}{2020--2024} \\
\cmidrule(lr){2-3}\cmidrule(lr){4-5}\cmidrule(lr){6-7}\cmidrule(l){8-9}
Spec. & Gross SR & Net SR & Gross SR & Net SR & Gross SR & Net SR & Gross SR & Net SR \\
\midrule
$I5/R5$   & 7.23 & 4.41 & 3.26 & 0.97 & 1.53 & $-0.15$ & 0.25 & $-0.86$ \\
$I5/R20$  & 1.81 & 0.83 & 0.31 & $-0.24$ & 0.11 & $-0.36$ & 0.39 & 0.07 \\
$I5/R60$  & 0.74 & 0.46 & 0.27 & $-0.04$ & 0.08 & $-0.07$ & $-0.25$ & $-0.39$ \\
$I20/R5$  & 7.76 & 5.12 & 3.19 & 1.11 & 1.74 & $-0.01$ & 0.90 & $-0.35$ \\
$I20/R20$ & 2.51 & 1.68 & 1.56 & 1.04 & 0.41 & $-0.02$ & 0.02 & $-0.38$ \\
$I20/R60$ & 0.60 & 0.40 & 0.38 & 0.25 & 0.05 & $-0.12$ & 0.19 & 0.08 \\
$I60/R5$  & 7.58 & 5.31 & 1.91 & 0.44 & 2.23 & 0.44 & 0.64 & $-0.75$ \\
$I60/R20$ & 1.89 & 1.31 & $-0.14$ & $-0.40$ & 0.75 & 0.35 & 0.02 & $-0.35$ \\
$I60/R60$ & 0.75 & 0.57 & 0.53 & 0.43 & 0.43 & 0.29 & 0.64 & 0.50 \\
\bottomrule
\end{tabular}
\begin{minipage}{\textwidth}
\footnotesize\textit{Note.} Every cell uses one selected feature-concatenation checkpoint. Each observation pairs chart and sequence inputs for the same stock, formation date, input length, and forecast horizon.
\end{minipage}
\end{table}

\FloatBarrier

\section{Rank stability across input windows}\label{app:rank_stability}

Table~\ref{tab:app_rank_stability} reports the underlying correlations in a compact tabular form. Each row holds the model family and return horizon fixed and compares the cross-sectional score rankings produced by two input lengths.

\begin{table}[htbp]
\centering
\caption{Within-model Spearman rank correlations across input windows}
\label{tab:app_rank_stability}
\small
\setlength{\tabcolsep}{9pt}
\begin{tabular}{@{}llrrr@{}}
\toprule
Model & Return horizon & $I5/I20$ & $I5/I60$ & $I20/I60$ \\
\midrule
Chart CNN & $R5$  & 0.60 & 0.22 & 0.34 \\
Chart CNN & $R20$ & 0.28 & 0.07 & 0.09 \\
Chart CNN & $R60$ & 0.18 & 0.09 & 0.33 \\
\addlinespace[2pt]
CNN1D & $R5$  & 0.64 & 0.54 & 0.71 \\
CNN1D & $R20$ & 0.52 & 0.38 & 0.62 \\
CNN1D & $R60$ & 0.30 & 0.21 & 0.45 \\
\addlinespace[2pt]
Multimodal & $R5$  & 0.56 & 0.44 & 0.59 \\
Multimodal & $R20$ & 0.23 & 0.19 & 0.37 \\
Multimodal & $R60$ & 0.21 & 0.11 & 0.32 \\
\bottomrule
\end{tabular}
\begin{minipage}{\textwidth}
\footnotesize\textit{Note.} Entries are time-series averages of formation-date Spearman correlations between two input-window scores from the same learned-model family over 2001--2024. The correlations are computed on the common stock-date panel used in the canonical overlap analysis.
\end{minipage}
\end{table}

\FloatBarrier

\section{Seed and fusion sensitivity}\label{app:sensitivity}

The chart seed comparison is available for the two short-input weekly specifications. Table~\ref{tab:app_ensemble_comparison} shows that the matched one- and five-member results are close in these cells. Table~\ref{tab:app_late_fusion} compares probability and logit averaging using the same aligned branch predictions.

\begin{table}[htbp]
\centering
\caption{Matched one- and five-member chart CNN results}
\label{tab:app_ensemble_comparison}
\scriptsize
\setlength{\tabcolsep}{4.3pt}
\begin{tabular}{@{}llrrrrrr@{}}
\toprule
Spec. & Weight & Gross SR, $M=1$ & Gross SR, $M=5$ & Net return, $M=1$ & Net return, $M=5$ & Net SR, $M=1$ & Net SR, $M=5$ \\
\midrule
$I5/R5$  & EW & 5.88 & 5.99 & 48.8\% & 49.9\% & 3.44 & 3.54 \\
$I5/R5$  & VW & 1.39 & 1.44 & $-2.1$\% & $-1.2$\% & $-0.13$ & $-0.07$ \\
$I20/R5$ & EW & 6.52 & 6.58 & 51.3\% & 55.3\% & 3.94 & 4.10 \\
$I20/R5$ & VW & 1.44 & 1.40 & $-2.5$\% & $-2.2$\% & $-0.16$ & $-0.13$ \\
\bottomrule
\end{tabular}
\begin{minipage}{\textwidth}
\footnotesize\textit{Note.} The table reports full-sample 2001--2024 portfolios under matched construction rules. $M$ is the number of independently trained chart forecasts averaged before the cross-sectional sort.
\end{minipage}
\end{table}

\begin{table}[htbp]
\centering
\caption{Late-fusion sensitivity to probability and logit averaging}
\label{tab:app_late_fusion}
\scriptsize
\setlength{\tabcolsep}{4.2pt}
\begin{tabular}{@{}llrrrrrr@{}}
\toprule
Spec. & Fusion rule & EW full & EW baseline & EW extension & VW full & VW baseline & VW extension \\
\midrule
$I5/R5$  & Probability average & 6.25 & 7.26 & 3.47 & 1.21 & 1.41 & 0.55 \\
$I5/R5$  & Logit average       & 6.25 & 7.25 & 3.47 & 1.20 & 1.40 & 0.55 \\
$I20/R5$ & Probability average & 6.51 & 7.70 & 3.40 & 1.60 & 1.89 & 0.82 \\
$I20/R5$ & Logit average       & 6.51 & 7.69 & 3.39 & 1.59 & 1.86 & 0.84 \\
\bottomrule
\end{tabular}
\begin{minipage}{\textwidth}
\footnotesize\textit{Note.} Entries are annualized gross Sharpe ratios. Baseline and extension denote 2001--2019 and 2020--2024. Late fusion combines the aligned predictions of pretrained chart and sequence branches.
\end{minipage}
\end{table}

\FloatBarrier

\section{Price-based benchmark grids}\label{app:benchmark_grids}

The benchmark tables use the same portfolio frequencies, weighting rules, annualization, and transaction-cost schedule as the learned portfolios. CONT is the implemented medium-horizon continuation signal. It compounds returns over trading-day lags 22 through 273, so its name records the implemented 252-day window without imposing the conventional 12-minus-1 label. STR and WSTR are one-month and one-week reversal signals. TREND is the supervised rolling linear forecast built from moving-average-to-price ratios.

\begin{table}[htbp]
\centering
\caption{Equal-weighted price-based benchmark Sharpe ratios}
\label{tab:app_benchmark_ew}
\scriptsize
\setlength{\tabcolsep}{4.6pt}
\begin{tabular}{@{}llrrrrrr@{}}
\toprule
& & \multicolumn{2}{c}{2001--2019} & \multicolumn{2}{c}{2020--2024} & \multicolumn{2}{c}{2001--2024} \\
\cmidrule(lr){3-4}\cmidrule(lr){5-6}\cmidrule(l){7-8}
Frequency & Signal & Gross SR & Net SR & Gross SR & Net SR & Gross SR & Net SR \\
\midrule
Week & CONT  & $-0.19$ & $-0.36$ & 0.26 & 0.12 & $-0.08$ & $-0.25$ \\
Week & STR   & 2.04 & 1.33 & 1.10 & 0.44 & 1.83 & 1.13 \\
Week & WSTR  & 3.33 & 1.59 & 1.42 & $-0.03$ & 2.86 & 1.20 \\
Week & TREND & 3.09 & 1.85 & 0.38 & $-0.62$ & 2.49 & 1.31 \\
\addlinespace[2pt]
Month & CONT  & $-0.01$ & $-0.09$ & 0.34 & 0.27 & 0.06 & $-0.01$ \\
Month & STR   & 0.71 & 0.32 & 0.24 & $-0.07$ & 0.59 & 0.22 \\
Month & WSTR  & 1.43 & 0.87 & 0.47 & 0.13 & 1.12 & 0.63 \\
Month & TREND & 0.70 & 0.37 & 0.60 & 0.30 & 0.68 & 0.35 \\
\addlinespace[2pt]
Quarter & CONT  & $-0.13$ & $-0.16$ & 0.38 & 0.35 & $-0.04$ & $-0.07$ \\
Quarter & STR   & 0.34 & 0.22 & 0.09 & $-0.00$ & 0.28 & 0.17 \\
Quarter & WSTR  & 0.37 & 0.18 & 0.56 & 0.30 & 0.39 & 0.19 \\
Quarter & TREND & $-0.05$ & $-0.11$ & 0.54 & 0.43 & 0.01 & $-0.06$ \\
\bottomrule
\end{tabular}
\begin{minipage}{\textwidth}
\footnotesize\textit{Note.} Gross and net Sharpe ratios are annualized. Net returns apply the common transaction-cost schedule to full absolute traded notional.
\end{minipage}
\end{table}

\begin{table}[htbp]
\centering
\caption{Value-weighted price-based benchmark Sharpe ratios}
\label{tab:app_benchmark_vw}
\scriptsize
\setlength{\tabcolsep}{4.6pt}
\begin{tabular}{@{}llrrrrrr@{}}
\toprule
& & \multicolumn{2}{c}{2001--2019} & \multicolumn{2}{c}{2020--2024} & \multicolumn{2}{c}{2001--2024} \\
\cmidrule(lr){3-4}\cmidrule(lr){5-6}\cmidrule(l){7-8}
Frequency & Signal & Gross SR & Net SR & Gross SR & Net SR & Gross SR & Net SR \\
\midrule
Week & CONT  & 0.12 & $-0.05$ & 0.53 & 0.38 & 0.21 & 0.05 \\
Week & STR   & 0.59 & 0.01 & $-0.51$ & $-1.01$ & 0.34 & $-0.23$ \\
Week & WSTR  & 0.82 & $-0.27$ & $-0.59$ & $-1.41$ & 0.47 & $-0.55$ \\
Week & TREND & 0.72 & $-0.09$ & $-0.21$ & $-0.77$ & 0.52 & $-0.24$ \\
\addlinespace[2pt]
Month & CONT  & 0.18 & 0.11 & 0.26 & 0.21 & 0.20 & 0.13 \\
Month & STR   & 0.24 & $-0.05$ & $-0.31$ & $-0.49$ & 0.09 & $-0.17$ \\
Month & WSTR  & 0.32 & $-0.01$ & $-0.25$ & $-0.46$ & 0.16 & $-0.13$ \\
Month & TREND & $-0.11$ & $-0.33$ & 0.39 & 0.24 & 0.02 & $-0.19$ \\
\addlinespace[2pt]
Quarter & CONT  & 0.07 & 0.03 & 0.52 & 0.48 & 0.16 & 0.12 \\
Quarter & STR   & 0.20 & 0.10 & 0.04 & $-0.02$ & 0.15 & 0.07 \\
Quarter & WSTR  & $-0.02$ & $-0.16$ & $-0.07$ & $-0.14$ & $-0.04$ & $-0.15$ \\
Quarter & TREND & $-0.03$ & $-0.08$ & 1.21 & 1.13 & 0.13 & 0.07 \\
\bottomrule
\end{tabular}
\begin{minipage}{\textwidth}
\footnotesize\textit{Note.} Gross and net Sharpe ratios are annualized. Value weights use lagged market capitalization and are normalized separately within the long and short legs.
\end{minipage}
\end{table}

\FloatBarrier

\section{Decile profiles}\label{app:decile_profiles}

The complete decile profiles show whether the high-minus-low result is supported by a broader cross-sectional ordering. The six panels report annualized gross decile returns for each model, input length, forecast horizon, and weighting rule over 2001--2024. They reproduce the return panels from the finalized simple-table exhibits and use no graphical interpolation. The full-grid Sharpe ratios appear in Tables~\ref{tab:app_chart_grid}--\ref{tab:app_mmd_grid}.

\begin{landscape}
\begin{table}[p]
\centering
\scriptsize
\setlength{\tabcolsep}{5pt}
\renewcommand{\arraystretch}{0.92}
\caption{Weekly learned-model decile performance, equal-weighted portfolios}
\label{tab:results_weekly_learned_deciles_ew}
\begin{tabular*}{\linewidth}{@{\extracolsep{\fill}}lrrrrrrrrr@{}}
\toprule
\multicolumn{10}{@{}l}{\textit{Annualized gross returns}} \\
\midrule
 & \multicolumn{3}{c}{Chart CNN} & \multicolumn{3}{c}{CNN1D} & \multicolumn{3}{c}{Multimodal} \\
\cmidrule(lr){2-4}\cmidrule(lr){5-7}\cmidrule(l){8-10}
Decile & $I5/R5$ & $I20/R5$ & $I60/R5$ & $I5/R5$ & $I20/R5$ & $I60/R5$ & $I5/R5$ & $I20/R5$ & $I60/R5$ \\
\midrule
Low & -27.1\% & -31.9\% & -16.4\% & -26.5\% & -34.6\% & -38.7\% & -27.0\% & -34.3\% & -35.2\% \\
2 & -3.8\% & -2.8\% & 2.2\% & -3.5\% & -4.7\% & -3.7\% & -5.0\% & -4.4\% & -4.0\% \\
3 & 2.9\% & 4.8\% & 9.4\% & 1.7\% & 4.9\% & 5.9\% & 1.8\% & 4.5\% & 6.9\% \\
4 & 6.9\% & 9.6\% & 12.5\% & 5.5\% & 9.5\% & 11.4\% & 6.7\% & 9.0\% & 10.9\% \\
5 & 10.3\% & 12.4\% & 15.0\% & 9.5\% & 13.6\% & 15.8\% & 10.9\% & 12.9\% & 14.6\% \\
6 & 13.8\% & 15.6\% & 16.8\% & 12.4\% & 16.6\% & 19.2\% & 13.2\% & 16.9\% & 17.8\% \\
7 & 18.1\% & 19.4\% & 18.8\% & 17.5\% & 20.1\% & 21.4\% & 18.6\% & 20.6\% & 21.0\% \\
8 & 23.5\% & 23.0\% & 21.7\% & 25.0\% & 23.7\% & 25.1\% & 24.7\% & 23.3\% & 24.3\% \\
9 & 32.1\% & 29.5\% & 23.4\% & 34.8\% & 29.4\% & 29.1\% & 31.8\% & 31.1\% & 29.4\% \\
High & 56.1\% & 53.0\% & 29.7\% & 56.4\% & 54.5\% & 49.5\% & 57.0\% & 53.3\% & 49.3\% \\
H--L & 83.3\% & 84.9\% & 46.1\% & 82.9\% & 89.1\% & 88.2\% & 84.0\% & 87.5\% & 84.5\% \\
\bottomrule
\end{tabular*}
\vspace{0.35em}
\parbox{\linewidth}{\scriptsize \textit{Note.} The table reports annualized gross decile returns over the available weekly observations in 2001--2024. Low is the lowest-score decile, High is the highest-score decile, and H--L is the long-short spread.}
\vspace{1.0em}
\scriptsize
\setlength{\tabcolsep}{5pt}
\renewcommand{\arraystretch}{0.92}
\caption{Weekly learned-model decile performance, value-weighted portfolios}
\label{tab:results_weekly_learned_deciles_vw}
\begin{tabular*}{\linewidth}{@{\extracolsep{\fill}}lrrrrrrrrr@{}}
\toprule
\multicolumn{10}{@{}l}{\textit{Annualized gross returns}} \\
\midrule
 & \multicolumn{3}{c}{Chart CNN} & \multicolumn{3}{c}{CNN1D} & \multicolumn{3}{c}{Multimodal} \\
\cmidrule(lr){2-4}\cmidrule(lr){5-7}\cmidrule(l){8-10}
Decile & $I5/R5$ & $I20/R5$ & $I60/R5$ & $I5/R5$ & $I20/R5$ & $I60/R5$ & $I5/R5$ & $I20/R5$ & $I60/R5$ \\
\midrule
Low & 0.1\% & -2.7\% & -0.8\% & 1.9\% & -0.2\% & -11.9\% & 1.9\% & -3.4\% & -5.6\% \\
2 & 5.3\% & 9.2\% & 4.8\% & 5.3\% & 4.5\% & -0.0\% & 3.0\% & 3.0\% & 5.4\% \\
3 & 4.8\% & 9.0\% & 6.6\% & 3.2\% & 6.5\% & 5.0\% & 8.0\% & 7.5\% & 5.0\% \\
4 & 6.6\% & 6.6\% & 8.7\% & 7.1\% & 7.6\% & 8.0\% & 7.0\% & 8.4\% & 6.1\% \\
5 & 9.4\% & 8.1\% & 6.9\% & 7.9\% & 7.6\% & 6.6\% & 9.0\% & 7.2\% & 8.3\% \\
6 & 9.5\% & 7.8\% & 11.6\% & 10.3\% & 10.5\% & 7.9\% & 9.6\% & 9.8\% & 8.2\% \\
7 & 10.0\% & 9.9\% & 9.9\% & 10.7\% & 10.5\% & 8.5\% & 11.5\% & 10.7\% & 9.1\% \\
8 & 13.3\% & 13.3\% & 11.2\% & 17.4\% & 11.2\% & 12.4\% & 12.3\% & 11.7\% & 12.0\% \\
9 & 14.6\% & 14.1\% & 12.4\% & 18.1\% & 15.6\% & 13.2\% & 16.9\% & 14.9\% & 13.1\% \\
High & 23.9\% & 20.5\% & 13.9\% & 23.2\% & 17.9\% & 20.0\% & 21.8\% & 20.2\% & 19.3\% \\
H--L & 23.8\% & 23.2\% & 14.7\% & 21.3\% & 18.0\% & 31.9\% & 19.9\% & 23.6\% & 25.0\% \\
\bottomrule
\end{tabular*}
\vspace{0.35em}
\parbox{\linewidth}{\scriptsize \textit{Note.} The table reports annualized gross decile returns over the available weekly observations in 2001--2024. Low is the lowest-score decile, High is the highest-score decile, and H--L is the long-short spread.}
\end{table}
\end{landscape}
\begin{landscape}
\begin{table}[p]
\centering
\scriptsize
\setlength{\tabcolsep}{5pt}
\renewcommand{\arraystretch}{0.92}
\caption{Monthly learned-model decile performance, equal-weighted portfolios}
\label{tab:results_monthly_learned_deciles_ew}
\begin{tabular*}{\linewidth}{@{\extracolsep{\fill}}lrrrrrrrrr@{}}
\toprule
\multicolumn{10}{@{}l}{\textit{Annualized gross returns}} \\
\midrule
 & \multicolumn{3}{c}{Chart CNN} & \multicolumn{3}{c}{CNN1D} & \multicolumn{3}{c}{Multimodal} \\
\cmidrule(lr){2-4}\cmidrule(lr){5-7}\cmidrule(l){8-10}
Decile & $I5/R20$ & $I20/R20$ & $I60/R20$ & $I5/R20$ & $I20/R20$ & $I60/R20$ & $I5/R20$ & $I20/R20$ & $I60/R20$ \\
\midrule
Low & -3.2\% & -2.1\% & 0.7\% & -9.2\% & -5.8\% & -9.7\% & -3.2\% & -3.1\% & -8.4\% \\
2 & 0.3\% & 5.1\% & 2.2\% & -2.2\% & 4.4\% & 4.2\% & -0.3\% & 5.6\% & -1.0\% \\
3 & 2.2\% & 7.7\% & 2.9\% & -0.3\% & 9.6\% & 7.8\% & 0.6\% & 9.2\% & 1.2\% \\
4 & 1.7\% & 10.2\% & 2.5\% & 1.5\% & 10.5\% & 11.5\% & -0.4\% & 10.1\% & 2.8\% \\
5 & 3.8\% & 12.7\% & 2.8\% & 0.9\% & 12.5\% & 13.8\% & 0.9\% & 11.2\% & 3.9\% \\
6 & 3.2\% & 12.1\% & 2.5\% & 1.6\% & 13.6\% & 16.1\% & 1.4\% & 13.4\% & 5.8\% \\
7 & 3.3\% & 14.2\% & 4.1\% & 4.4\% & 13.6\% & 17.1\% & 2.7\% & 14.1\% & 6.6\% \\
8 & 4.8\% & 13.8\% & 4.1\% & 4.6\% & 16.0\% & 17.0\% & 1.8\% & 14.9\% & 6.2\% \\
9 & 5.4\% & 17.7\% & 4.2\% & 5.6\% & 17.2\% & 17.8\% & 5.8\% & 16.4\% & 7.7\% \\
High & 10.6\% & 21.8\% & 6.7\% & 11.9\% & 22.0\% & 20.2\% & 9.5\% & 21.8\% & 10.3\% \\
H--L & 13.8\% & 24.0\% & 6.0\% & 21.0\% & 27.8\% & 29.8\% & 12.7\% & 24.9\% & 18.7\% \\
\bottomrule
\end{tabular*}
\vspace{0.35em}
\parbox{\linewidth}{\scriptsize \textit{Note.} The table reports annualized gross decile returns over the available monthly observations in 2001--2024. Low is the lowest-score decile, High is the highest-score decile, and H--L is the long-short spread.}
\vspace{1.0em}
\scriptsize
\setlength{\tabcolsep}{5pt}
\renewcommand{\arraystretch}{0.92}
\caption{Monthly learned-model decile performance, value-weighted portfolios}
\label{tab:results_monthly_learned_deciles_vw}
\begin{tabular*}{\linewidth}{@{\extracolsep{\fill}}lrrrrrrrrr@{}}
\toprule
\multicolumn{10}{@{}l}{\textit{Annualized gross returns}} \\
\midrule
 & \multicolumn{3}{c}{Chart CNN} & \multicolumn{3}{c}{CNN1D} & \multicolumn{3}{c}{Multimodal} \\
\cmidrule(lr){2-4}\cmidrule(lr){5-7}\cmidrule(l){8-10}
Decile & $I5/R20$ & $I20/R20$ & $I60/R20$ & $I5/R20$ & $I20/R20$ & $I60/R20$ & $I5/R20$ & $I20/R20$ & $I60/R20$ \\
\midrule
Low & 1.2\% & 6.4\% & 4.5\% & 1.1\% & 5.5\% & 4.0\% & 3.0\% & 6.1\% & -5.1\% \\
2 & 0.6\% & 5.4\% & 2.7\% & 0.2\% & 8.3\% & 4.6\% & 1.0\% & 6.5\% & -0.8\% \\
3 & 1.2\% & 7.6\% & 1.2\% & -0.2\% & 8.7\% & 8.6\% & 1.4\% & 10.6\% & -3.3\% \\
4 & -1.1\% & 7.3\% & 3.8\% & 1.5\% & 8.8\% & 6.7\% & 2.9\% & 9.1\% & 3.4\% \\
5 & 3.0\% & 7.8\% & 3.4\% & 1.7\% & 9.7\% & 9.1\% & 0.2\% & 9.2\% & 1.2\% \\
6 & 0.1\% & 8.7\% & -1.0\% & 2.9\% & 10.5\% & 9.3\% & 1.8\% & 10.6\% & 1.4\% \\
7 & 4.3\% & 10.8\% & 2.0\% & 0.9\% & 8.7\% & 12.0\% & 3.6\% & 9.6\% & 3.4\% \\
8 & 3.1\% & 8.8\% & 0.4\% & 4.1\% & 8.9\% & 10.2\% & 2.2\% & 8.9\% & 2.0\% \\
9 & 1.5\% & 12.3\% & -0.3\% & 3.4\% & 11.1\% & 10.6\% & 1.7\% & 10.1\% & 3.0\% \\
High & 2.6\% & 12.4\% & 1.0\% & 7.0\% & 11.6\% & 12.4\% & 5.0\% & 10.5\% & 4.0\% \\
H--L & 1.4\% & 6.0\% & -3.5\% & 5.9\% & 6.1\% & 8.4\% & 2.0\% & 4.4\% & 9.1\% \\
\bottomrule
\end{tabular*}
\vspace{0.35em}
\parbox{\linewidth}{\scriptsize \textit{Note.} The table reports annualized gross decile returns over the available monthly observations in 2001--2024. Low is the lowest-score decile, High is the highest-score decile, and H--L is the long-short spread.}
\end{table}
\end{landscape}
\begin{landscape}
\begin{table}[p]
\centering
\scriptsize
\setlength{\tabcolsep}{5pt}
\renewcommand{\arraystretch}{0.92}
\caption{Quarterly learned-model decile performance, equal-weighted portfolios}
\label{tab:results_quarterly_learned_deciles_ew}
\begin{tabular*}{\linewidth}{@{\extracolsep{\fill}}lrrrrrrrrr@{}}
\toprule
\multicolumn{10}{@{}l}{\textit{Annualized gross returns}} \\
\midrule
 & \multicolumn{3}{c}{Chart CNN} & \multicolumn{3}{c}{CNN1D} & \multicolumn{3}{c}{Multimodal} \\
\cmidrule(lr){2-4}\cmidrule(lr){5-7}\cmidrule(l){8-10}
Decile & $I5/R60$ & $I20/R60$ & $I60/R60$ & $I5/R60$ & $I20/R60$ & $I60/R60$ & $I5/R60$ & $I20/R60$ & $I60/R60$ \\
\midrule
Low & 6.6\% & 5.0\% & 1.4\% & 7.0\% & 8.0\% & 7.3\% & 5.9\% & 1.9\% & 0.1\% \\
2 & 6.2\% & 5.5\% & 5.0\% & 10.9\% & 10.7\% & 10.2\% & 6.4\% & 5.7\% & 3.1\% \\
3 & 7.2\% & 6.4\% & 7.2\% & 11.0\% & 12.0\% & 11.1\% & 4.7\% & 7.9\% & 6.7\% \\
4 & 6.2\% & 7.3\% & 8.0\% & 11.6\% & 12.6\% & 12.5\% & 6.5\% & 7.0\% & 8.1\% \\
5 & 8.5\% & 7.5\% & 8.4\% & 13.1\% & 13.7\% & 13.4\% & 6.3\% & 7.9\% & 9.5\% \\
6 & 8.3\% & 9.0\% & 8.7\% & 13.2\% & 12.7\% & 14.6\% & 7.6\% & 9.0\% & 9.2\% \\
7 & 7.8\% & 7.9\% & 8.6\% & 13.8\% & 14.5\% & 14.4\% & 9.4\% & 8.0\% & 9.7\% \\
8 & 8.8\% & 9.0\% & 10.5\% & 13.6\% & 13.4\% & 15.1\% & 9.2\% & 10.0\% & 10.7\% \\
9 & 9.0\% & 10.7\% & 9.9\% & 16.0\% & 15.5\% & 16.2\% & 9.9\% & 10.7\% & 11.8\% \\
High & 9.1\% & 9.8\% & 11.1\% & 18.2\% & 15.7\% & 17.1\% & 11.7\% & 10.0\% & 12.3\% \\
H--L & 2.5\% & 4.8\% & 9.7\% & 11.2\% & 7.7\% & 9.8\% & 5.8\% & 8.1\% & 12.2\% \\
\bottomrule
\end{tabular*}
\vspace{0.35em}
\parbox{\linewidth}{\scriptsize \textit{Note.} The table reports annualized gross decile returns over the available quarterly observations in 2001--2024. Low is the lowest-score decile, High is the highest-score decile, and H--L is the long-short spread.}
\vspace{1.0em}
\scriptsize
\setlength{\tabcolsep}{5pt}
\renewcommand{\arraystretch}{0.92}
\caption{Quarterly learned-model decile performance, value-weighted portfolios}
\label{tab:results_quarterly_learned_deciles_vw}
\begin{tabular*}{\linewidth}{@{\extracolsep{\fill}}lrrrrrrrrr@{}}
\toprule
\multicolumn{10}{@{}l}{\textit{Annualized gross returns}} \\
\midrule
 & \multicolumn{3}{c}{Chart CNN} & \multicolumn{3}{c}{CNN1D} & \multicolumn{3}{c}{Multimodal} \\
\cmidrule(lr){2-4}\cmidrule(lr){5-7}\cmidrule(l){8-10}
Decile & $I5/R60$ & $I20/R60$ & $I60/R60$ & $I5/R60$ & $I20/R60$ & $I60/R60$ & $I5/R60$ & $I20/R60$ & $I60/R60$ \\
\midrule
Low & 11.0\% & 5.4\% & 2.8\% & 7.7\% & 10.5\% & 9.9\% & 9.6\% & 7.8\% & 1.8\% \\
2 & 9.4\% & 4.7\% & 6.1\% & 10.4\% & 8.9\% & 10.4\% & 9.9\% & 3.3\% & 4.4\% \\
3 & 8.6\% & 7.1\% & 8.5\% & 9.1\% & 10.3\% & 10.6\% & 5.3\% & 9.0\% & 6.5\% \\
4 & 9.3\% & 11.4\% & 6.8\% & 12.1\% & 12.8\% & 10.0\% & 6.2\% & 7.8\% & 8.5\% \\
5 & 4.8\% & 6.8\% & 6.0\% & 12.5\% & 11.1\% & 10.0\% & 5.6\% & 9.0\% & 6.3\% \\
6 & 7.8\% & 5.8\% & 6.2\% & 10.4\% & 9.8\% & 10.9\% & 6.6\% & 9.0\% & 8.3\% \\
7 & 9.0\% & 10.6\% & 8.0\% & 10.9\% & 10.5\% & 10.7\% & 8.4\% & 7.4\% & 5.4\% \\
8 & 6.5\% & 8.6\% & 8.1\% & 10.9\% & 9.6\% & 11.1\% & 9.5\% & 7.0\% & 9.3\% \\
9 & 8.8\% & 6.7\% & 8.5\% & 9.9\% & 10.8\% & 10.5\% & 10.7\% & 7.7\% & 11.0\% \\
High & 8.5\% & 7.8\% & 9.0\% & 11.5\% & 12.6\% & 13.8\% & 8.9\% & 7.8\% & 8.8\% \\
H--L & -2.5\% & 2.4\% & 6.2\% & 3.7\% & 2.1\% & 3.9\% & -0.7\% & 0.1\% & 7.0\% \\
\bottomrule
\end{tabular*}
\vspace{0.35em}
\parbox{\linewidth}{\scriptsize \textit{Note.} The table reports annualized gross decile returns over the available quarterly observations in 2001--2024. Low is the lowest-score decile, High is the highest-score decile, and H--L is the long-short spread.}
\end{table}
\end{landscape}


\FloatBarrier
